\chapter*{Conclusion générale et perspectives}
\addcontentsline{toc}{chapter}{Conclusion générale et perspectives}
\label{chap:conclusion}

Le projet \textit{VitalSync} a abouti à la conception et au développement d’un prototype innovant dédié à la surveillance en temps réel des signes vitaux en salle de réveil post-opératoire, répondant aux besoins spécifiques des hôpitaux tunisiens, souvent confrontés à des ressources limitées. Réalisé au sein du SmartLab de la Faculté de Médecine de Monastir, ce projet a été mené par une équipe interdisciplinaire selon la méthodologie Scrum, en quatre sprints, et a mobilisé des technologies avancées telles que la reconnaissance optique de caractères (OCR), l’Internet des Objets (IoT) et les notifications en temps réel via Firebase Cloud Messaging (FCM). Ce chapitre récapitule les principales réalisations, évalue les retombées du système et propose des pistes pour son évolution future.

Le Sprint 0 a permis de poser les fondations du projet, en identifiant les besoins des utilisateurs finaux et en définissant une architecture modulaire adaptée, avec une répartition claire des rôles entre les acteurs (médecins, infirmiers) et les composants du système (base de données, API, interfaces).

Le Sprint 1 a marqué une avancée technique notable grâce à la mise en place d’un pipeline OCR basé sur OpenCV, YOLO et Tesseract. Ce pipeline a permis l’extraction automatisée des paramètres vitaux à partir de moniteurs hétérogènes, apportant une réponse concrète à la problématique de compatibilité multi-marques rencontrée dans les environnements hospitaliers tunisiens.

Lors du Sprint 2, le système a été enrichi d’un mécanisme de détection d’anomalies et de génération d’alertes, avec l’intégration de notifications mobiles via FCM et la création d’un tableau de bord initial pour les professionnels de santé. Cette phase a renforcé la dimension proactive du système, facilitant une réaction rapide face aux situations cliniques critiques.

Le Sprint 3 a finalisé le système \textit{VitalSync}, en optimisant ses performances et en améliorant l’expérience utilisateur, tant sur le plan fonctionnel qu’ergonomique. L’interface web et mobile a été enrichie de filtres dynamiques et de visualisations graphiques, consolidant ainsi l’outil de suivi à destination du personnel médical. L’architecture logicielle a été stabilisée, et l’environnement de production préparé, ouvrant la voie à un déploiement futur.

\textit{VitalSync} se distingue par son architecture modulaire, sa compatibilité avec plusieurs marques de moniteurs médicaux, et son approche économique, adaptée au contexte tunisien. Comparé aux solutions propriétaires existantes, il propose une automatisation accrue de la collecte des données, une diffusion des alertes en temps réel, et une intégration technique plus souple aux infrastructures hospitalières existantes. Il contribue ainsi à améliorer la réactivité du personnel médical, à alléger leur charge de travail et à renforcer la sécurité des patients.

Néanmoins, ce prototype nécessite encore des ajustements pour un déploiement à grande échelle. Parmi les perspectives envisagées figurent : la conduite de tests cliniques en conditions réelles, l’adaptation du système à des situations de terrain plus variées (pannes réseau, éclairage variable, etc.), ainsi que l’intégration avec des dossiers médicaux électroniques (DME) et des standards tels que HL7 pour favoriser l’interopérabilité. À plus long terme, l’intégration de mécanismes d’apprentissage automatique pourrait permettre de prédire les variations des signes vitaux et d’anticiper les situations critiques. Une collaboration renforcée avec les autorités de santé et des partenaires industriels serait également un levier important pour assurer le financement, la validation réglementaire et la diffusion de la solution à plus grande échelle.

En conclusion, \textit{VitalSync} constitue une avancée significative vers la modernisation des soins post-opératoires en Tunisie. En combinant une approche agile, des technologies modernes et une vision centrée sur les utilisateurs, le projet a démontré la faisabilité d’une solution de surveillance vitale temps réel, flexible et accessible. Son évolution future dépendra d’un effort coordonné entre acteurs académiques, professionnels de santé et partenaires techniques, dans l’objectif commun d’améliorer durablement la qualité des soins dans les milieux hospitaliers à ressources limitées.